\documentclass[11pt]{article}
\usepackage{fontspec}
\setmainfont{Times New Roman}
\setsansfont{Arial}
\setmonofont{Consolas}
\usepackage{amsmath,amssymb,mathtools}
\usepackage{geometry}
\usepackage{microtype}
\geometry{a4paper,margin=2.35cm}
\setlength{\parindent}{1.25em}
\setlength{\parskip}{0pt}
\makeatletter
\renewcommand\footnotesize{\@setfontsize\footnotesize{7.7}{8.7}\selectfont}
\makeatother
\emergencystretch=35em
\allowdisplaybreaks
\sloppy
\hbadness=10000
\vbadness=10000
\hfuzz=2pt
\vfuzz=10pt
\raggedbottom

\providecommand{\mR}{\mathfrak{R}}
\providecommand{\mA}{\mathfrak{A}}
\providecommand{\mB}{\mathfrak{B}}
\providecommand{\mC}{\mathfrak{C}}
\providecommand{\ma}{\mathfrak{a}}
\providecommand{\mb}{\mathfrak{b}}
\providecommand{\mc}{\mathfrak{c}}
\providecommand{\mm}{\mathfrak{m}}
\providecommand{\mo}{\mathfrak{o}}
\providecommand{\mv}{\mathfrak{v}}

\begin{document}

% Unit: Paper30 Section04 v001. Direct Interslavic Latin translation from the German final-audited source slice, source lines 542--596 / segments P30-S0065--P30-S0089. Footnote counter continues after Paper30 Section03.
\setcounter{footnote}{14}

\subsection*{\S~4. Teoremy o izomorfiji. Direktne sumy}

V slědujučem predpostavlja se samo kolcevo, odnosno modulno svojstvo, ale nikaka druga aksioma.

1. Ako $M$ i $\overline M$ sut modulne oblasti vzhodno k $\mR$ (§2), togda $M$ nazyvaje se homomorfny k $\overline M$ (točnije: modulno-homomorfny), $M\sim\overline M$, ako vsakomu elementu iz $M$ sootvětny jest jeden i samo jeden element iz $\overline M$ tako, že tym $\overline M$ jest vyčerpany;\footnote{Vse v smyslu definicije ravnosti, ktora vlada v $\overline M$; sr. primětku 6.} i ako pri toj sootvětnosti razlika i množenje tym samym elementom iz $\mR$ sootvětny sut jedno drugomu. To jest, iz $\beta\sim\overline\beta$ i $\gamma\sim\overline\gamma$ vsegda slěduje:
\[
(\beta-\gamma)\sim(\overline\beta-\overline\gamma),\qquad r\beta\sim r\overline\beta .
\]

Zato $\mR$-modulu $\mB$ iz $M$ homomorfno sootvětny jest $\mR$-modul $\overline{\mB}$ iz $\overline M$; a cělost $\mC^*$ elementov iz $M$, ktorym sootvětny sut elementy $\mR$-modula $\overline{\mC}$ iz $\overline M$, tvori v $M$ modul $\mC^*$, jednoznačno opreděljen črez $\overline{\mC}$, asociovany s $\overline{\mC}$, pri čem $\mC^*\sim\overline{\mC}$. Ako osoblivo $\mA$ jest modul, asociovany s nulovym elementom iz $\overline M$, togda $\mC^*$ jest dělitelj $\mA$ i razpada se na klasy elementov iz $M$, kongruentnyh modulo $\mA$, tako že te klasy vzaimno jednoznačno sootvětny sut elementam iz $\overline{\mC}$. Ako se prějde od $\mB$ v $M$ k $\overline{\mB}$ v $\overline M$, a potom nazad k $\mB^*$, dobiva se $\mB^*=(\mB,\mA)$, to jest največši obči dělitelj $\mB$ i $\mA$; zato $\mB^*=\mB$, ako $\mB$ jest dělitelj $\mA$.

Ako sootvětnost medžu elementami iz $M$ i $\overline M$ jest obratno jednoznačna, oblasti nazyvajut se izomorfne (točnije: modulno-izomorfne), $M\cong\overline M$.

Ako $\mA$ jest libovolny $\mR$-modul iz $M$, togda nastava modulna oblast $\overline M$, homomorfna k $M$ -- modul klasov ostatkov $M\mid\mA$ -- kogda kongruencija po modulu $\mA$ pojmaje se kako nova relacija ravnosti. Vsakomu elementu iz $M$ pri tom prirědžujut se vse i samo te elementy, ktore sut ravne jemu v $\overline M$. Prěhod v $\overline M$ od definicije ravnosti k identitetě znači, že vse ravne elementy v $\overline M$ sobirajut se v jednu klasu -- ostatkovu klasu -- i te ostatkove klasy pojmajut se kako nove elementy $\overline M$.

\emph{Vsaky homomorfizm nastava prěhodom k modulu klasov ostatkov;} bo ako $M\sim\overline M$, i ako $\mA$ jest modul asociovany s nulovym elementom iz $\overline M$, togda, kako pokazano vyše, $\overline M$ jest izomorfny k modulu klasov ostatkov $M\mid\mA$.

2. \textbf{Prvy teorem o izomorfiji.} \emph{Ako $\overline M$ označa modul klasov ostatkov $M\mid\mA$, i ako $\mC$ jest dělitelj $\mA$, togda važi izomorfizm
\[
\overline M\mid\overline{\mC}\cong M\mid\mC .
\]}
Bo kongruencija po modulu $\mA$ jest jednočasno kongruencija po modulu $\mC$; elementy ravne po modulu $\mA$ ostavajut ravne po modulu $\mC$. Zato možno tvoriti modul klasov ostatkov $M\mid\mC$ -- to jest ravnost po modulu $\mC$ -- tako, že se najprvo elementy izjednačujut po modulu $\mA$, tedy prěhodi se k $\overline M$, i medžu nimi sobirajut se te, ktore sut ravne po modulu $\mC$; v $\overline M$ to sootvětny jest izjednačenju po modulu $\overline{\mC}$, tedy tvorjenju $\overline M\mid\overline{\mC}$.

\textbf{Drugi teorem o izomorfiji.} \emph{Ako $\mB$ i $\mA$ sut moduly iz $M$, togda važi izomorfizm
\[
(\mB,\mA)\mid\mA\cong \mB\mid[\mB,\mA].
\]}
Bo po 1. $\mB$ stava homomorfny k $\overline{\mB}$, ako $(\mB,\mA)\mid\mA$ položi se ravnym $\overline{\mB}$; a poneže pri tom nulovomu elementu v $\overline{\mB}$ sootvětny sut vse elementy iz $[\mB,\mA]$ i samo oni, vyše dany izomorfizm znovo slěduje iz 1.

3. Ako $\mR$ i $\overline{\mR}$ sut (komutativne) kolca, togda $\mR$ nazyvaje se homomorfno k $\overline{\mR}$ (točnije: kolcevo-homomorfno), $\mR\sim\overline{\mR}$, ako vsakomu elementu iz $\mR$ sootvětny jest jeden i samo jeden element iz $\overline{\mR}$ tako, že tym $\overline{\mR}$ jest vyčerpano, i ako pri toj prirědbě razlike i produkty sootvětny sut jedno drugomu. Ako sootvětnost elementov jest obratno jednoznačna, kolca nazyvajut se izomorfne (točnije: kolcevo-izomorfne), $\mR\cong\overline{\mR}$.

Vse razsmatrjanja iz 1. i 2. ostavajut validne, ako moduly zaměniti idealami v $\mR$,\footnote{Pri nekomutativnyh kolcah treba vzęti za osnovu dvustranne idealy.} i ako pojem modulnogo homomorfizma i modulnogo izomorfizma zaměniti kolcevym homomorfizmom i kolcevym izomorfizmom. Osoblivo, ako $\mc$ jest dělitelj $\ma$, kolco klasov ostatkov $\mc\mid\ma$ nastava, kogda kongruencija po modulu $\ma$ uvodi se kako nova relacija ravnosti; pri tom treba iměti v vidu, že teper sut definirane takodže produkty ostatkovyh klasov.

Togda $\mc$ jest homomorfno k $\mc\mid\ma$; vsaky homomorfizm nastava na toj način, i teoremy o izomorfiji važe kako v 2:

\textbf{Prvy teorem o izomorfiji.} Ako $\overline{\mR}$ označa kolco klasov ostatkov $\mR\mid\ma$, i ako $\mc$ jest dělitelj $\ma$, togda $\overline{\mR}\mid\overline{\mc}\cong\mR\mid\mc$ v smyslu kolcevogo izomorfizma.

\textbf{Drugi teorem o izomorfiji.} Ako $\mb$ i $\ma$ sut idealy iz $\mR$, togda važi kolcevy izomorfizm:
\[
(\mb,\ma)\mid\ma\cong\mb\mid[\mb,\ma].
\]\footnote{Te teoremy o izomorfiji, znane iz teorije grup, najprvo pojavjajut se za moduly v trochu specialnějšej formě u Dedekinda (sr. 3. izd., str. 484). Vyše dana forma za idealy nahodi se u Masazo Sono: \emph{On Congruences. I, II, III, IV}, Memoirs of the College of Science, Kyoto Imperial University, 2 (1917), 3 (1918, 1919), sr. 2, str. 215. V vsakom slučaju izrično formulovany jest samo drugi teorem o izomorfiji.}

4. V slědujučem budut sobrané računske pravila za vzaimno proste idealy,\footnote{I to v formě, ktora idě nazad k Dedekindu (4. izd., §178, III--VIII). Računanje s jediničnym elementom, ktoro povsudu pojavjaje se v novějših izloženjah, mnogo jest složnějše.} iz ktoryh v 5. budut slědovati teoremy o direktnyh sumah. V komutativnom kolcu $\mR$ predpostavlja se eksistencija jediničnogo elementa $e$ množenja. Zato $\ma=\mo\ma$ za vsaky ideal iz $\mR$, kde pod $\mo$ razuměje se jediničny ideal. Dva idealy $\ma,\mb$ nazyvajut se vzaimno proste, ako jih največši obči dělitelj $(\ma,\mb)$ stava ravny $\mo$.

\emph{4$\alpha$. Ako $(\ma,\mb)=\mo$, togda $(\ma,\mb\mc)=(\ma,\mc)$.} Bo
\[
(\ma,\mc)=(\ma,\mb)\cdot(\ma,\mc)=(\ma^2,\ma\mb,\ma\mc,\mb\mc)
\]
jest dělimo črez $(\ma,\mb\mc)$ i obratno. Iz togo slěduje:

\emph{4$\beta$. Iz $\mc\mb\equiv0\pmod{\ma}$ i $(\ma,\mb)=\mo$ slěduje $\mc\equiv0\pmod{\ma}$.} Bo $\mc\mb\equiv0\pmod{\ma}$ jest ravnoznačno s $(\ma,\mb\mc)=\ma$; zato zaradi $(\ma,\mb)=\mo$ takodže $(\ma,\mc)=\ma$, ili $\mc\equiv0\pmod{\ma}$. Dalej slěduje:

\emph{4$\gamma$. Ako každy iz idealov $\ma_1,\ldots,\ma_r$ jest vzaimno prosty s každym iz idealov $\mb_1,\ldots,\mb_s$, togda produkt $\ma_i$ jest vzaimno prosty s produktom $\mb_i$.} Bo iz $(\ma,\mb)=\mo$ i $(\ma,\mc)=\mo$ dobiva se $(\ma,\mb\mc)=\mo$, a konečno povtorjenje daje tvrdženje.

Iz 4$\alpha$ i 4$\gamma$ dobiva se:

\emph{4$\delta$. Ako idealy $\mb_1,\ldots,\mb_r$ sut poparno vzaimno proste, to jest $(\mb_i,\mb_k)=\mo$ za $i\ne k$, togda jih najmenše obče mnogokratno jest ravno jih produktu.} Naj $(\ma,\mb)=\mo$, $\mv=[\ma,\mb]$. Togda
\[
\mv=\mo\mv=(\ma\mv,\mb\mv)=0(\ma,\mb);
\]
poněže $\ma\mb\equiv0\pmod{\mv}$, odtud slěduje $\mv=\ma\mb$, a konečno povtorjenje s użyvanjem 4$\gamma$ daje tvrdženje.

\emph{4$\varepsilon$. Ako idealy $\mb_1,\ldots,\mb_r$ sut poparno vzaimno proste, i ako}
\[
\ma_i=[\mb_1,\ldots,\mb_{i-1},\mb_{i+1},\ldots,\mb_r]
=\mb_1\cdots\mb_{i-1}\mb_{i+1}\cdots\mb_r
\]
\emph{jest položeno, togda $(\ma_1,\ldots,\ma_{i-1},\ma_{i+1},\ldots,\ma_r)=\mb_i$, i zato $(\ma_1,\ldots,\ma_r)=\mo$.} Bo po distributivnom zakonu
\[
(\ma_1,\ma_2)=\mb_3\cdots\mb_r(\mb_2,\mb_1)=\mb_3\cdots\mb_r,
\]
i zato
\[
(\ma_1,\ma_2,\ma_3)=\mb_4\cdots\mb_r(\mb_3,\mb_1\mb_2)=\mb_4\cdots\mb_r,
\]
iz čego konečnym povtorjenjem pri vhodnom numerovanju slěduje tvrdženje.

Ideal $\ma_i$ nazyvaje se komplementom $\mb_i$ v predstavjenju $\mm=[\mb_1,\ldots,\mb_r]$.

5. V komutativnom kolcu $\mR$ znovo predpostavlja se eksistencija jediničnogo elementa množenja. Kolco $\mR$ nazyvaje se direktna suma idealov $\ma_1,\ldots,\ma_r$ -- v znakih
\[
\mR=\ma_1+\ma_2+\cdots+\ma_r,
\]
ako vsaky element $c$ iz $\mR$ možno predstaviti jednym i samo jednym načinom v formě $c=a_1+a_2+\cdots+a_r$, kde vsaky $a_i$ jest element iz $\ma_i$.

\emph{Ako nulovy ideal jest najmenše obče mnogokratno poparno vzaimno prostyh idealov $\mb_1,\ldots,\mb_r$, i ako $\ma_i$ jest komplement $\mb_i$, togda $\mR$ jest direktna suma idealov $\ma_1,\ldots,\ma_r$.} Bo zaradi $\mo=(\ma_1,\ldots,\ma_r)$ vsaky element iz $\mR$ dopustja najmene jedno aditivno predstavjenje črez $\ma_i$; zaradi $[\ma_i,\mb_i]=(0)$ to predstavjenje jest jednoznačno, bo iz $0=a_1+a_2+\cdots+a_r=a_i+b_i$ dobiva se $a_i=0$. Po drugom teoremu o izomorfiji idealy $\ma_i$ sut izomorfne k kolcam klasov ostatkov $\mR\mid\mb_i$. Važe "ortogonalne relacije" $\ma_i\ma_k=(0)$ za $i\ne k$ i $\ma_i^2=\ma_i$; poslědnje zaradi $\ma_i=\mo\ma_i$.

\emph{Obratno, ako eksistuje predstavjenje $\mR$ kako direktnoj sumy,
\[
\mR=\ma_1+\ma_2+\cdots+\ma_r,
\]
i ako položi se
\[
\mb_i=(\ma_1,\ldots,\ma_{i-1},\ma_{i+1},\ldots,\ma_r)
=\ma_1+\cdots+\ma_{i-1}+\ma_{i+1}+\cdots+\ma_r,
\]
togda $\mb_i$ sut poparno vzaimno proste i jih najmenše obče mnogokratno jest ravno nulovomu idealu.} Bo iz $\mo=(\ma_i,\mb_i)$ slěduje $\mo=(\mb_k,\mb_i)$ za $k\ne i$, poněže $\mb_k$ jest dělitelj $\ma_i$. Naj dalje $c$ jest dělimo črez vse $\mb_i$. Togda
\[
 c=a_2^{(1)}+\cdots+a_r^{(1)}
 =a_1^{(2)}+a_3^{(2)}+\cdots+a_r^{(2)}
 =\cdots=a_1^{(r)}+\cdots+a_{r-1}^{(r)},
\]
kde vsakokrat $a_i^{(\lambda)}\equiv0\pmod{\ma_i}$. Iz jednoznačnosti predstavjenja, poněže vsakokrat jedna komponenta stava nulova, slěduje $0=a_1^{(\lambda)},\ldots,0=a_r^{(\lambda)}$ za vsako $\lambda$, i tym $c=0$.\footnote{Teoremy dane pod 5. sut specialne slučaji obćego teorema o vezi medžu direktnoj sumoju i direktnym prěsěkom. Sr. H. Prüfer, Theorie der Abelschen Gruppen I, Math. Zeitschr. 20 (1924), str. 165--187, §6. Tam dane razsmatrjanja ostavajut validne takodže pri tako obćem formulovanju pojma grupy, že moduly i idealy podredžujut se jemu kako specialne slučaji.}

\end{document}
